A questão que se nos apresenta é de uma beleza singular, pois trata de um dos operadores lineares
mais fundamentais e visualizáveis: a projeção. Através de sua análise, veremos como a estrutura de
decomposição de um espaço vetorial em soma direta reflete-se diretamente nas propriedades
espectrais (autovalores, autovetores) e nos polinômios associados ao operador.

Vamos dissecar o problema com a precisão que a Álgebra exige.

\subsection*{Enunciado do Problema:}

Seja $V$ um espaço vetorial tal que $V = W_1 \oplus W_2$, com $\dim(W_1)=n$ e $\dim(W_2)=m$. Seja
$T:V \to V$ o operador linear de projeção de $V$ sobre $W_1$ ao longo de $W_2$. Determine os
autovalores, os autovetores, o polinômio característico e o polinômio minimal de $T$.

\subsection*{Resolução Detalhada}

A chave para desvendar este problema reside na correta interpretação da definição de projeção no
contexto de uma soma direta.

A hipótese $V=W_1 \oplus W_2$ significa que todo vetor $v \in V$ pode ser escrito, de maneira \textit{única},
como a soma $v=w_1+w_2$, onde $w_1 \in W_1$ e $w_2 \in W_2$.

O operador $T$ é a projeção sobre $W_1$ ao longo de $W_2$. Por definição, sua ação sobre um vetor $v$ é
dada por:

\[
T(v) = T(w_1+w_2)=w_1
\]

Com esta ferramenta fundamental em mãos, podemos proceder à análise.

\subsubsection*{1. Autovalores e Autovetores}

Por definição, um escalar $\lambda$ é um autovalor de $T$ se existe um vetor não nulo $v$ (o autovetor) tal que
$T(v)=\lambda v$. Vamos investigar a ação de $T$ sobre os vetores dos subespaços $W_1$ e $W_2$.

\textbf{Caso 1: Análise sobre o subespaço $W_1$}

Seja $v$ um vetor não nulo pertencente a $W_1$, isto é, $v \in W_1$ e $v \neq 0$.
Para aplicar o operador $T$, devemos decompor $v$ na forma $w_1+w_2$. Como $v \in W_1$, sua
decomposição única é:

\[
v = v+0
\]

onde $v \in W_1$ e $0 \in W_2$ (o vetor nulo).

Aplicando $T$:

\[
T(v)=T(v+0)=v
\]

Comparando com a equação de autovalor, $T(v)=\lambda v$, temos:

\[
v=\lambda v \;\;\Longrightarrow\;\; 1 \cdot v = \lambda v
\]

Como $v\neq 0$, concluímos que $\lambda=1$.

Isto significa que \textbf{todo vetor não nulo de $W_1$ é um autovetor de $T$ associado ao autovalor $\lambda=1$}.
O autoespaço associado a $\lambda=1$ é, portanto, o próprio subespaço $W_1$.

\textbf{Caso 2: Análise sobre o subespaço $W_2$}

Seja $v$ um vetor não nulo pertencente a $W_2$, isto é, $v \in W_2$ e $v \neq 0$.
Sua decomposição única na soma direta é:

\[
v = 0+v
\]

onde $0 \in W_1$ e $v \in W_2$.

Aplicando $T$:

\[
T(v) = T(0+v)=0
\]

Comparando com a equação de autovalor, $T(v)=\lambda v$, temos:

\[
0=\lambda v
\]

Como $v\neq 0$, a única solução para esta equação é $\lambda=0$.

Isto significa que \textbf{todo vetor não nulo de $W_2$ é um autovetor de $T$ associado ao autovalor $\lambda=0$}.
O autoespaço associado a $\lambda=0$ é o próprio subespaço $W_2$, que corresponde ao núcleo do
operador, $\ker(T)$.

\textbf{Conclusão sobre Autovalores e Autovetores:}

Os autovalores de $T$ são $\lambda_1=1$ e $\lambda_2=0$.

\begin{itemize}
\item O autoespaço associado a $\lambda=1$ é $V_1=W_1$.
\item O autoespaço associado a $\lambda=0$ é $V_0=W_2$.
\end{itemize}

Como $V=W_1 \oplus W_2$, o espaço $V$ é a soma direta dos autoespaços de $T$. Isso implica que o
operador $T$ é diagonalizável.

\subsubsection*{2. Polinômio Característico}

Para determinar o polinômio característico $p_T(t)=\det(T-tI)$, a estratégia mais elegante é
construir uma base para $V$ que seja adaptada à decomposição em soma direta.

Seja $\mathcal{B}_1=\{u_1,u_2,\ldots,u_n\}$ uma base para $W_1$.

Seja $\mathcal{B}_2=\{v_1,v_2,\ldots,v_m\}$ uma base para $W_2$.

Como $V=W_1 \oplus W_2$, a união ordenada $\mathcal{B}=\{u_1,\ldots,u_n,v_1,\ldots,v_m\}$ forma uma base para $V$.
A dimensão de $V$ é $n+m$.

Vamos encontrar a representação matricial de $T$ nesta base, $[T]_{\mathcal{B}}$:

\begin{itemize}
\item Para os vetores de $\mathcal{B}_1$: $T(u_i)=u_i=1\cdot u_i$. As coordenadas de $T(u_i)$ na base $\mathcal{B}$ são
$(0,\ldots,1,\ldots,0)^T$, com $1$ na $i$-ésima posição.
\item Para os vetores de $\mathcal{B}_2$: $T(v_j)=0$. As coordenadas de $T(v_j)$ na base $\mathcal{B}$ são o vetor nulo
$(0,\ldots,0)^T$.
\end{itemize}

A matriz de $T$ na base $\mathcal{B}$ é uma matriz em blocos:

\[
[T]_{\mathcal{B}} =
\begin{pmatrix}
I_n & 0 \\
0 & 0_m
\end{pmatrix}
\]

onde $I_n$ é a matriz identidade $n\times n$ e $0_m$ é a matriz nula $m\times m$.

Agora, calculamos o polinômio característico:

\[
p_T(t) = \det([T]_{\mathcal{B}} - tI_{n+m})
\]

O determinante de uma matriz diagonal em blocos é o produto dos determinantes dos blocos:

\[
p_T(t) = \det((1-t)I_n)\cdot \det(-tI_m)
\]

\[
p_T(t) = (1-t)^n \cdot (-t)^m
\]

Portanto, o polinômio característico de $T$ é $p_T(t)=(-1)^m t^m (1-t)^n$.
As raízes são $t=0$ com multiplicidade algébrica $m=\dim(W_2)$, e $t=1$ com multiplicidade
algébrica $n=\dim(W_1)$, como esperado.

\subsubsection*{3. Polinômio Minimal}

O polinômio minimal, $m_T(t)$, é o polinômio mônico (coeficiente do termo de maior grau igual a $1$) de
menor grau que anula o operador $T$, i.e., $m_T(T)=0$.

Vamos analisar a natureza do operador de projeção. O que acontece se aplicarmos $T$ duas vezes?
Seja $v=w_1+w_2$ um vetor qualquer em $V$.

\[
T(v)=w_1
\]

Como $w_1 \in W_1$, sua decomposição é $w_1=w_1+0$. Logo:

\[
T(w_1)=w_1
\]

Assim, para todo $v \in V$, temos $T^2(v)=w_1=T(v)$.
Isso nos mostra que $T^2=T$, ou seja, $T$ é um operador idempotente.

A relação $T^2=T$ pode ser reescrita como $T^2-T=0$. Isso significa que o polinômio $q(t)=t^2-t=t(t-1)$ é um polinômio aniquilador de $T$.

Pelo Teorema do Polinômio Minimal, sabemos que $m_T(t)$ deve dividir qualquer polinômio aniquilador
de $T$. Logo, $m_T(t)$ deve dividir $q(t)=t(t-1)$.
Os divisores mônicos de $t(t-1)$ são: $1, t, t-1$ e $t(t-1)$.

\begin{itemize}
\item Se $m_T(t)=t$, então $T=0$. Isso só ocorre se $W_1=\{0\}$ (i.e., $n=0$), um caso trivial.
\item Se $m_T(t)=t-1$, então $T-I=0$, ou $T=I$. Isso só ocorre se $W_2=\{0\}$ (i.e., $m=0$), outro caso trivial.
\end{itemize}

Assumindo o caso não-trivial onde $n>0$ e $m>0$, nem $T$ nem $T-I$ são o operador nulo.
Portanto, os polinômios de grau 1 não podem aniquilar $T$.

Concluímos que o polinômio mônico de menor grau que anula $T$ é:

\[
m_T(t)=t(t-1)=t^2-t
\]

(Nos casos triviais, se $n=0$, $T=0$ e $m_T(t)=t$. Se $m=0$, $T=I$ e $m_T(t)=t-1$).

\subsection*{Sumário dos Resultados}

Para o operador de projeção $T$ sobre $W_1$ ao longo de $W_2$, com $\dim(W_1)=n>0$ e
$\dim(W_2)=m>0$:

\begin{itemize}
\item \textbf{Autovalores:} $\lambda_1=1$ e $\lambda_2=0$.
\item \textbf{Autovetores:} O autoespaço de $\lambda=1$ é $W_1$. O autoespaço de $\lambda=0$ é $W_2$.
\item \textbf{Polinômio Característico:} $p_T(t)=(-1)^m t^m (1-t)^n$.
\item \textbf{Polinômio Minimal:} $m_T(t)=t(t-1)$.
\end{itemize}

O fato de o polinômio minimal se fatorar em termos lineares distintos é uma confirmação algébrica
da nossa conclusão geométrica de que o operador $T$ é diagonalizável.